% Generated draft: requires mathematical and source review.
\documentclass[11pt]{article}
\usepackage[margin=25mm]{geometry}
\usepackage{fontspec}
\setmainfont{Latin Modern Roman}
\usepackage{amsmath,amssymb,mathtools,unicode-math}
\setmathfont{Latin Modern Math}
\usepackage{enumitem,needspace,xurl,hyperref}
\setlength{\parindent}{0pt}
\setlength{\parskip}{5pt}
\setlength{\emergencystretch}{3em}
\newcommand{\N}{\mathbb{N}}
\newcommand{\Z}{\mathbb{Z}}
\newcommand{\Q}{\mathbb{Q}}
\newcommand{\R}{\mathbb{R}}
\newcommand{\C}{\mathbb{C}}
\newcommand{\F}{\mathbb{F}}
\newcommand{\A}{\mathbb{A}}
\newcommand{\PP}{\mathbb{P}}
\newcommand{\E}{\mathbb{E}}
\newcommand{\eps}{\varepsilon}
\newcommand{\dd}{\,\mathrm{d}}
\begin{document}
\section*{UnsolvedMath ID: 700100}
\textbf{Source problem code:} catalogue entry
\subsection{Definitions and mathematical statement}
Let \(G\) be a group with identity \(1\). For \(x,y\in G\), define the iterated commutators by
\[
[x,{}_1y]=x^{-1}y^{-1}xy,\qquad [x,{}_{n+1}y]=[[x,{}_ny],y]\quad(n\geq 1).
\]
The group \(G\) is an Engel group if
\[
\forall x,y\in G\ \exists n\in\Z_{\geq 1}\quad [x,{}_ny]=1.
\]
Here \(n\) may depend on both \(x\) and \(y\); no uniform bound is assumed. The group is torsion-free if
\[
\forall g\in G\ \forall m\in\Z_{\geq 1}\quad (g^m=1\Longrightarrow g=1).
\]
A strict total order \(<\) on \(G\) is left-invariant if
\[
\forall a,b,c\in G\quad (a<b\Longrightarrow ca<cb),
\]
and is bi-invariant if
\[
\forall a,b,c,d\in G\quad (a<b\Longrightarrow cad<cbd).
\]
A group admitting such an order is called left-orderable or bi-orderable, respectively.

\textbf{Source question (A. I. Kokorin, 1967).} Is every torsion-free Engel group orderable? No finiteness, finite-generation, or countability hypothesis is supplied.

\textbf{Unresolved qualification.} The supplied record does not specify the invariance meant by ``orderable.'' Thus the intended target cannot be uniquely fixed from this record. The two corresponding precise formulations are
\[
\begin{aligned}
&\forall G\quad (G\text{ is torsion-free and Engel}\Longrightarrow G\text{ is left-orderable}),\\
&\forall G\quad (G\text{ is torsion-free and Engel}\Longrightarrow G\text{ is bi-orderable}).
\end{aligned}
\]
These are alternative interpretations of the source question, not two separately attributed conjectures. The intended convention requires source review.
\subsection{Short English statement}
A. I. Kokorin asks whether every torsion-free Engel group admits a compatible total order. The Engel condition allows the commutator length to depend on the chosen pair of elements. The supplied record leaves unclear whether the order must be left-invariant or bi-invariant.
\subsection{Sources}
\begin{itemize}
\item Supplied local UnsolvedMath registry record: permanent ID \(700100\), exact source problem code \texttt{catalogue entry}, title ``Kourovka Notebook Problem 2.24.'' Upstream revision: \texttt{d371ee917ac9647fbb06cd6fc529deec6f3dda82}.
\item \emph{The Kourovka Notebook (Unsolved Problems in Group Theory)}, original problem 2.24, 2nd issue (1967), proposed by A. I. Kokorin; identified in the supplied background with the 21st issue edition (2026), \texttt{arXiv:1401.0300v48}. The supplied collection description names E. I. Khukhro and V. D. Mazurov as editors.
\item The registry labels this problem ``open'' and records its update at \texttt{2026-10-16T14:00:00Z}. This is supplied registry information, not an independently verified mathematical or literature assessment. No separate literature-check date or literature-source links are supplied.
\end{itemize}
\end{document}
